%% toc: P15 165 165.6
\chapter{Partition Updates, Evidence Models and Independence for Lower Previsions}
\label{sec:p15-w02b-p15-034-partition-updates}

\Cref{sec:p15-ip-conditioning} conditions a lower prevision on one
observation event at a time. Learning which cell of a partition obtains is
a different operation: one conditional lower prevision is required for every
cell simultaneously, and the whole family has to be coherent with the
unconditional assessment. For a finite partition the single-event results
can be applied cell by cell. For an infinite partition they cannot, and
this chapter shows exactly why. It declares the partition-update semantics
used in this book, proves that such an update exists if and only if the
unconditional model has the conglomerability property for that partition,
identifies the least update when it exists, and works a finitely additive
precise prevision for which every cell has positive probability but no
coherent partition update exists. The development follows the definitions
of Miranda's author manuscript
\cite[author manuscript, §§3.1--3.3, pp.~17--24]{MirandaCoherentLowerPrevisionsCapture};
the proofs written in the annex are the ones relied upon here.
Two further questions are settled in the same framework. The first is how
lower previsions, belief functions and possibility measures differ once
gambles and updates are taken into account, and not only event values. The
second is which independence judgement a decision between two acts on a
pair of variables actually needs when the marginal models are imprecise.

\emph{Inputs.} A nonempty possibility space \(\Omega\), a partition of it,
and a coherent lower prevision defined on every bounded gamble on
\(\Omega\). \emph{Conclusions.} Whether the input admits any conditional
lower prevision on the partition that is coherent with it; if so, the
least such one and the values every such one must take; if not, an explicit
finite-stake family of conditional and unconditional transactions with a
uniformly negative net gain. For the evidence models of
\cref{sec:p15-w02b-p15-034-evidence-models} the inputs are a finite
mass assignment or possibility distribution, and the conclusions are its
gamble extension and the comparison of Dempster's conditioning with the
coherent update. For \cref{sec:p15-w02b-p15-034-independence-decision} the
inputs are two imprecise binary marginals, a stated independence judgement,
two acts and a comparison criterion; the conclusion is whether the model
ranks the acts, and under which judgement.

\section{Partitions, separate coherence and joint coherence}
\label{sec:p15-w02b-p15-034-definitions}

Throughout this chapter the event algebra of \cref{def:p15-ip-domain} is
the power set of \(\Omega\), so the gamble space is
\(\mathcal L(\Omega)\), the set of all bounded real functions on
\(\Omega\). The letter \(B\) therefore names a partition cell below,
following the source, and not the gamble space \(B(\Sigma)\) of
\cref{def:p15-ip-domain}. The unconditional input
\(\underline P\) is a coherent lower prevision on all of
\(\mathcal L(\Omega)\); by \cref{thm:p15-ip-natural-extension} it equals
its own natural extension, so the properties listed there and in
\cref{thm:p15-ip-linear-domain} hold for it. As in
\cref{def:p15-ip-coherence}, \(G(f)=f-\underline P(f)\), and
\(\overline P(f)=-\underline P(-f)\).

\begin{hbdefinition}{def:p15-w02b-p15-034-partition}{Partition and separately coherent conditional lower prevision}
A \emph{partition} \(\mathcal B\) of \(\Omega\) is a set of nonempty,
pairwise disjoint subsets of \(\Omega\) whose union is \(\Omega\); it may
be finite, countably infinite or uncountable. For \(\omega\in\Omega\)
write \(B_\omega\) for the cell containing \(\omega\).

A \emph{conditional lower prevision on \(\mathcal B\)} assigns to every
cell \(B\in\mathcal B\) a real functional
\(\underline P(\cdot\mid B)\) on \(\mathcal L(\Omega)\). It is
\emph{separately coherent} when for every \(B\in\mathcal B\), every
\(f,g\in\mathcal L(\Omega)\) and every real \(\lambda>0\),
\[
\begin{aligned}
 &\text{(C1)}\quad \underline P(f\mid B)\geq\inf_{\omega\in B}f(\omega),\\
 &\text{(C2)}\quad \underline P(\lambda f\mid B)=\lambda\underline P(f\mid B),\\
 &\text{(C3)}\quad \underline P(f+g\mid B)\geq\underline P(f\mid B)+\underline P(g\mid B).
\end{aligned}
\]
These are the three conditions the source gives for a domain that is a
linear space containing the constants \cite[p.~18]{MirandaCoherentLowerPrevisionsCapture};
the definition here takes them as the requirement.
For \(f\in\mathcal L(\Omega)\) define the gambles
\[
\begin{gathered}
 \underline P(f\mid\mathcal B)(\omega)=\underline P(f\mid B_\omega),\qquad
 G(f\mid B)=I_B\bigl(f-\underline P(f\mid B)\bigr),\\
 G(f\mid\mathcal B)=f-\underline P(f\mid\mathcal B)=\sum_{B\in\mathcal B}G(f\mid B),
\end{gathered}
\]
the sum being pointwise with at most one nonzero term at each point
\cite[eq.~(10), p.~18]{MirandaCoherentLowerPrevisionsCapture}.
The \(\mathcal B\)-\emph{support} of \(g\) is
\(S(g)=\{B\in\mathcal B:I_Bg\ne0\}\). For an unconditional gamble
\(f\), put \(S_\Omega(f)=\{\Omega\}\) if \(f\ne0\) and
\(S_\Omega(f)=\varnothing\) if \(f=0\).
\end{hbdefinition}

The value \(\underline P(f\mid B)\) is read as a present supremum buying
price for \(f\) called off unless the outcome lies in \(B\). Reading it as
the price the subject would pay after learning only that the outcome lies in
\(B\) adds Walley's \emph{updating principle}: a gamble is acceptable
contingent on \(B\) exactly when its product with \(I_B\) is acceptable now
\cite[§3.1, p.~17]{MirandaCoherentLowerPrevisionsCapture}. Joint coherence
for a partition rests on a second premise, the \emph{conglomerative
principle}: a gamble acceptable contingent on every cell of
\(\mathcal B\) is acceptable outright
\cite[§3.2, p.~19]{MirandaCoherentLowerPrevisionsCapture}. Both are
modelling commitments. Neither follows from the numerical bounds.

\begin{hbdefinition}{def:p15-w02b-p15-034-joint}{Joint coherence of an unconditional and a partition model}
Let \(\underline P\) be coherent on \(\mathcal L(\Omega)\) and let
\(\underline P(\cdot\mid\mathcal B)\) be separately coherent. They are
\emph{jointly coherent} when, for all
\(f_1,f_2,g_1,g_2\in\mathcal L(\Omega)\) and every \(B_0\in\mathcal B\),
\[
 \text{(JC1)}\quad
 \sup_{\omega\in\Omega}\bigl[G(f_1)+G(g_1\mid\mathcal B)-G(f_2)\bigr](\omega)\geq0,
\]
and (JC2): with
\(Z=G(f_1)+G(g_1\mid\mathcal B)-G(g_2\mid B_0)\), there is a set
\(C\in\{B_0\}\cup S(g_1)\cup S_\Omega(f_1)\) such that
\(\sup_{\omega\in C}Z(\omega)\geq0\).

These are the two conditions (JC1) and (JC2) of the source
\cite[§3.3, p.~22]{MirandaCoherentLowerPrevisionsCapture}, with the
supremum taken over the sets prescribed by its general coherence
condition~(11) \cite[§3.2, p.~19]{MirandaCoherentLowerPrevisionsCapture},
in which the unconditional model is the conditional model on the one-cell
partition \(\{\Omega\}\). In (JC1) the sold gamble is unconditional, so
\(\Omega\) belongs to the prescribed family and the supremum is over
\(\Omega\).
\end{hbdefinition}

The test sets matter. Taking every supremum over all of \(\Omega\) instead
is Walley's weaker notion, which the source warns is satisfied
automatically by any conditional portfolio whose active cells do not cover
\(\Omega\) \cite[Remark~1, p.~20]{MirandaCoherentLowerPrevisionsCapture}.
A portfolio in (JC1) or (JC2) may involve infinitely many cells at once,
through \(G(g_1\mid\mathcal B)\); this is where infinite partitions differ
from a single event.

\begin{hbdefinition}{def:p15-w02b-p15-034-conglomerable}{Conglomerability for a partition}
A coherent lower prevision \(\underline P\) on \(\mathcal L(\Omega)\)
is \emph{\(\mathcal B\)-conglomerable} when, for every sequence
\(B_1,B_2,\ldots\) of distinct cells of \(\mathcal B\) with
\(\underline P(B_n)>0\) for all \(n\), and every
\(f\in\mathcal L(\Omega)\) with \(\underline P(I_{B_n}f)\geq0\) for all
\(n\),
\[
             \underline P\bigl(I_{\bigcup_nB_n}f\bigr)\geq0
\]
\cite[§3.3, p.~23]{MirandaCoherentLowerPrevisionsCapture}.
It is \emph{fully conglomerable} when this holds for every partition
\cite[p.~24]{MirandaCoherentLowerPrevisionsCapture}.
\end{hbdefinition}

\begin{hbconvention}{def:p15-w02b-p15-034-semantics}{Declared semantics of a partition update}
In this book a \emph{partition update} of a coherent lower prevision
\(\underline P\) on \(\mathcal L(\Omega)\) relative to a partition
\(\mathcal B\) is a separately coherent conditional lower prevision
\(\underline P(\cdot\mid\mathcal B)\) on \(\mathcal L(\Omega)\) that is
jointly coherent with \(\underline P\) in the sense of
\cref{def:p15-w02b-p15-034-joint}. Its value \(\underline P(f\mid B)\) is
the supremum buying price for \(f\) contingent on \(B\); it is a price
after learning \(B\) only under the updating principle, and the joint test
across cells presupposes the conglomerative principle.

When a partition update exists, the \emph{least partition update}
\(\underline P_{\mathcal B}\) is
\[
 \underline P_{\mathcal B}(f\mid B)=
 \begin{cases}
  \text{the unique }q\in\mathbb R\text{ with }\underline P(I_B[f-q])=0,
                    &\underline P(B)>0,\\[2pt]
  \inf_{\omega\in B}f(\omega), &\underline P(B)=0.
 \end{cases}
\]
The first line is the generalized Bayes rule; the second is the vacuous
conditional model. A user who wants a more informative value on a cell of
lower probability zero must state that further assessment; it is not
supplied by \(\underline P\). On a cell of positive lower probability
the first line is the value that \cref{thm:p15-ip-generalized-bayes}
identifies with regular conditioning on that single event; a test that only
checks finitely many cells at once is a different criterion
(\cref{thm:p15-w02b-p15-034-williams}) and is not identified with this one.
\end{hbconvention}

\section{What joint coherence forces, and when an update exists}
\label{sec:p15-w02b-p15-034-results}

\begin{hblemma}{thm:p15-w02b-p15-034-cell-facts}{proved=annex:p15.w02b-p15-034.cell-facts}{Cellwise facts}
Let \(\underline P\) be coherent on \(\mathcal L(\Omega)\) and
\(\underline P(\cdot\mid\mathcal B)\) separately coherent.
\begin{enumerate}
\item For each cell \(B\),
\(\inf_Bf\leq\underline P(f\mid B)\leq\sup_Bf\), and
\(\underline P(f\mid B)=\underline P(g\mid B)\) whenever \(f=g\) on
\(B\). In particular \(\underline P(f\mid\mathcal B)\) is a gamble, and
\(G(I_Bf\mid\mathcal B)=G(f\mid B)\).
\item For each cell \(B\) and \(f,g\), \(\sup_B\bigl[(f-\underline P(f\mid B))-(g-\underline P(g\mid B))\bigr]\geq0\).
\item For a cell \(B\) and \(f\), the function
\(F(q)=\underline P(I_B[f-q])\) is \(1\)-Lipschitz and nonincreasing, and
\(F(q+\delta)\leq F(q)-\delta\underline P(B)\) for \(\delta>0\).
If \(\underline P(B)>0\) it has exactly one root, and that root lies in
\([\inf_Bf,\sup_Bf]\).
\item The cells with \(\underline P(B)>0\) form a finite or countably
infinite set; for each \(m\geq1\) at most \(m\) cells satisfy
\(\underline P(B)\geq1/m\).
\item If only finitely many cells have positive lower probability, every
finite family \(B_1,\ldots,B_k\) of them and every \(f\) with
\(\underline P(I_{B_i}f)\geq0\) satisfy
\(\underline P(I_{B_1\cup\cdots\cup B_k}f)\geq0\), and \(\underline P\)
is \(\mathcal B\)-conglomerable. This includes every finite partition and
every partition all of whose cells have lower probability zero.
\end{enumerate}
\end{hblemma}

\begin{hbtheorem}{thm:p15-w02b-p15-034-necessity}{proved=annex:p15.w02b-p15-034.necessity}{Joint coherence forces the generalized Bayes rule and conglomerability}
Let \(\underline P(\cdot\mid\mathcal B)\) be a partition update of
\(\underline P\). Then:
\begin{enumerate}
\item \(\underline P(G(f\mid B))=0\) for every \(f\in\mathcal L(\Omega)\)
and every \(B\in\mathcal B\);
\item if \(\underline P(B)>0\), \(\underline P(f\mid B)\) is the unique
root of \(q\mapsto\underline P(I_B[f-q])\); if \(\underline P(B)=0\),
\(\underline P(f\mid B)\geq\inf_Bf\);
\item \(\underline P\) is \(\mathcal B\)-conglomerable.
\end{enumerate}
\end{hbtheorem}

\begin{hbtheorem}{thm:p15-w02b-p15-034-existence}{proved=annex:p15.w02b-p15-034.existence}{Existence and the least partition update}
A coherent lower prevision \(\underline P\) on \(\mathcal L(\Omega)\) has
a partition update relative to \(\mathcal B\) if and only if it is
\(\mathcal B\)-conglomerable. In that case \(\underline P_{\mathcal B}\)
of \cref{def:p15-w02b-p15-034-semantics} is well defined and is a partition
update, and every partition update
\(\underline P(\cdot\mid\mathcal B)\) satisfies
\[
 \underline P(f\mid B)=\underline P_{\mathcal B}(f\mid B)\ \text{if }\underline P(B)>0,\qquad
 \underline P(f\mid B)\geq\underline P_{\mathcal B}(f\mid B)\ \text{if }\underline P(B)=0 .
\]
\end{hbtheorem}

By \cref{thm:p15-w02b-p15-034-cell-facts}, a finite partition always has
the least update, computed cell by cell exactly as in
\cref{sec:p15-ip-conditioning}. The content of
\cref{thm:p15-w02b-p15-034-existence} lies with partitions having
infinitely many cells of positive lower probability. The source states the
same equivalence \cite[p.~23]{MirandaCoherentLowerPrevisionsCapture}, and
also that the vacuous model is the least coherent choice on a cell of
lower probability zero \cite[p.~24]{MirandaCoherentLowerPrevisionsCapture};
both are proved here for the full-domain case.

\section{A finitely additive prevision with no partition update}
\label{sec:p15-w02b-p15-034-counterexample}

The construction needs a finitely additive probability on the positive
integers that gives every finite set probability zero. Its existence uses
the representation theorem, hence the Choice-dependent separation premise
already assumed in \cref{sec:p15-imprecise-probability}.

\begin{hbtheorem}{thm:p15-w02b-p15-034-counterexample}{proved=annex:p15.w02b-p15-034.counterexample}{Coherent and cellwise updatable, but not updatable on the partition}
Let \(\mathbb N_+=\{1,2,3,\ldots\}\),
\(\Omega=\mathbb N_+\times\{0,1\}\), and let
\(\mathcal B=\{B_n:n\in\mathbb N_+\}\) with
\(B_n=\{(n,0),(n,1)\}\). Let \(Y_1=\mathbb N_+\times\{1\}\).
\begin{enumerate}
\item Some positive normalized linear functional \(\nu\)
on \(\mathcal L(\mathbb N_+)\) has \(\nu(I_F)=0\) for every finite
\(F\subseteq\mathbb N_+\).
\item Define, for \(f\in\mathcal L(\Omega)\),
\[
\begin{gathered}
 \mu(f)=\sum_{n=1}^\infty2^{-n}\,\frac{f(n,0)+f(n,1)}2,\qquad
 \nu_1(f)=\nu\bigl(n\mapsto f(n,1)\bigr),\\
 P=\tfrac12(\mu+\nu_1).
\end{gathered}
\]
Then \(P\) is a linear prevision, hence a coherent lower prevision on
\(\mathcal L(\Omega)\) equal to its conjugate upper prevision. It satisfies
\(P(B_n)=2^{-n-1}\) for every \(n\), so
\(\sum_nP(B_n)=1/2\) although the cells cover \(\Omega\);
\(P(Y_1)=3/4\); and \(P(I_{B_n}Y_1)/P(B_n)=1/2\) for every \(n\).
\item For \(f=I_{\Omega\setminus Y_1}-I_{Y_1}\),
\(P(I_{B_n}f)=0\) for every \(n\) while \(P(f)=-1/2\). Hence \(P\) is
not \(\mathcal B\)-conglomerable, and \(P\) has no partition update
relative to \(\mathcal B\).
\item For each cell, the cell average
\(a_n(g)=\tfrac12[g(n,0)+g(n,1)]\) satisfies (C1)--(C3) on \(B_n\),
and \(a_n(g)\) is the unique root of
\(q\mapsto P(I_{B_n}[g-q])\). Using all these values together,
the portfolio that buys \(f\) contingent on every cell at price \(0\)
and sells \(f\) unconditionally at price \(-1/2\) has net gain
\(G(f\mid\mathcal B)-G(f)\) identically equal to \(-1/2\).
\item By contrast \(\mu\) alone is \(\mathcal B\)-conglomerable, and its
least partition update is \(a_n\) on \(B_n\), with
\(\mu(Y_1)=1/2\) equal to every cell value.
\end{enumerate}
\end{hbtheorem}

Every number in item~2 is fixed. Unconditionally the event \(Y_1\) has
probability \(3/4\); conditionally on whichever cell obtains it has
probability \(1/2\). Learning the cell is certain to lower the
probability, which is the defect the conglomerative principle excludes. The
cause is the purely finitely additive component \(\nu_1\): it gives each
cell probability zero but \(Y_1\) probability one, so half of the total
mass sits on no cell. Single-event conditioning sees nothing wrong, because
every individual cell has positive probability and its update is the
ordinary ratio; the incoherence appears only in a portfolio active on
infinitely many cells. Unconditional coherence of \(P\) is therefore
insufficient to justify the simultaneous conditional claims
\(P(Y_1\mid B_n)=1/2\) for all \(n\), which are the only values the
generalized Bayes rule allows cell by cell.

\begin{hbtheorem}{thm:p15-w02b-p15-034-williams}{proved=annex:p15.w02b-p15-034.williams}{The finite-support criterion accepts the same assessments}
In \cref{thm:p15-w02b-p15-034-counterexample}, assign the conditioning
event \(\Omega\) the prevision \(P\) and each cell \(B_n\) the prevision
\(a_n\). For every \(k\geq0\), every choice of conditioning events
\(C_0,\ldots,C_k\in\{\Omega\}\cup\mathcal B\), every
\(\lambda_0,\ldots,\lambda_k\geq0\) and every
\(f_0,\ldots,f_k\in\mathcal L(\Omega)\),
\[
 \sup_{\omega\in C_0\cup\cdots\cup C_k}
 \Bigl[\sum_{i=1}^k\lambda_iI_{C_i}\bigl(f_i-P(f_i\mid C_i)\bigr)
       -\lambda_0I_{C_0}\bigl(f_0-P(f_0\mid C_0)\bigr)\Bigr](\omega)\geq0,
\]
where \(P(\cdot\mid\Omega)=P\) and \(P(\cdot\mid B_n)=a_n\). This is
Williams' condition~(13) in the source's statement
\cite[§5.2, p.~32]{MirandaCoherentLowerPrevisionsCapture}.
The uniformly losing portfolio of
\cref{thm:p15-w02b-p15-034-counterexample} uses all cells at once and is
not of this form.
\end{hbtheorem}

The source describes condition~(13) as the case of the general coherence
condition~(11) in which each support is finite, states that for finite
partitions the two criteria agree, and notes that the finite-support
criterion does not guarantee conglomerability
\cite[§5.2, p.~32]{MirandaCoherentLowerPrevisionsCapture}.
\Cref{thm:p15-w02b-p15-034-williams} exhibits that difference on one
explicit model: the cellwise assessments pass every finitely supported test,
and fail the partition test of \cref{def:p15-w02b-p15-034-joint}.

\begin{hbwarning}{def:p15-w02b-p15-034-limits}{What these results do not settle}
The results above concern one partition and one unconditional model
defined on every gamble. They do not prove anything about several
partitions at once, about restricted assessment domains, or about the
natural extension of conditional assessments on infinite partitions. For
those cases the source reports, citing Walley's monograph, that a coherent
extension may fail to exist, that the least coherent extension may differ
from the natural extension, and that coherent conditional models need not
be lower envelopes of precise ones
\cite[§3.2.1, pp.~20--21]{MirandaCoherentLowerPrevisionsCapture}; those
reports are not used in this book's derivations. The vacuous choice on cells
of lower probability zero is the least one by
\cref{thm:p15-w02b-p15-034-existence}; whether a particular more
informative choice on infinitely many such cells, for instance regular
extension in each cell, is jointly coherent is not decided here and must
be checked against (JC1) and (JC2) for the model at hand.
Finally, accepting conglomerability is a substantive position. De Finetti
rejected it, and the source ties full conglomerability of linear
previsions to countable additivity
\cite[pp.~24 and~31]{MirandaCoherentLowerPrevisionsCapture}; compare the
countable additivity condition of \cref{thm:p15-ip-countable-additivity}.
\end{hbwarning}

\section{Belief functions, possibility and lower previsions}
\label{sec:p15-w02b-p15-034-evidence-models}

This section works on a finite nonempty space \(W\), with the mass
assignments \(m\), belief and plausibility functions
\(\operatorname{Bel}_m,\operatorname{Pl}_m\), possibility distributions
\(\pi\), possibility and necessity measures
\(\operatorname{Poss}_\pi,\operatorname{Nec}_\pi\), and allocation sets
\(\mathcal C_m\) of \cref{def:p14-evidence-uncertainty-data}. On a finite
space the gamble space of \cref{def:p15-ip-domain} is all of
\(\mathbb R^W\), and a precise model of \cref{def:p15-ip-compatible} is a
probability mass \(p\) with expectation \(E_p\). An event function
\(\nu\) on \(2^W\) is used as a lower probability by taking the
assessment domain to be the event indicators and
\(\underline P(I_A)=\nu(A)\); its compatible set is then
\[
 \mathcal M(\nu)=\{p:\ P_p(A)\geq\nu(A)\text{ for every }A\subseteq W\}.
\]
\Cref{thm:p14-evidence-finite-measures,thm:p14-evidence-measure-separation}
compare the event values of these constructions. The results below add the
two points at which they part company as models of belief: the values they
give to gambles, and the updates they license. The source surveyed here
treats belief functions as the completely monotone lower probabilities on a
finite space, reports that 2-monotone lower probabilities are coherent, and
concludes that the evidence theory embeds in the theory of coherent lower
previsions when updating and combination are left aside
\cite[§5.5, p.~34]{MirandaCoherentLowerPrevisionsCapture}. Those two
operations are where the theories part, and the results below make the
difference exact.

\begin{hbtheorem}{thm:p15-w02b-p15-034-belief-extension}{proved=annex:p15.w02b-p15-034.bel-ext}{The gamble extension of a belief function}
Let \(m\) be a mass assignment on \(W\).
\begin{enumerate}
\item As a lower probability on all events, \(\operatorname{Bel}_m\) is
coherent, and \(\mathcal C_m\subseteq\mathcal M(\operatorname{Bel}_m)\).
\item Its natural extension to all gambles is
\[
 \underline E_m(g)=\sum_{B\subseteq W,\ m(B)>0}m(B)\min_{w\in B}g(w)
 =\min_{p\in\mathcal C_m}E_p(g)\qquad(g\in\mathbb R^W).
\]
\item \(\underline E_m(g)\) equals the Choquet value
\(V_{\operatorname{Bel}_m}(g)\) of \cref{def:p17-rat84-choquet}, with
\(g\) in place of the utility vector.
\end{enumerate}
\end{hbtheorem}

A belief function therefore fixes a single gamble model, its natural
extension, and that extension is computed by sending each mass to the
worst point of its focal set. This is a commitment of the belief-function
reading. A coherent lower prevision whose values on events happen to form a
belief function need not make it, as the next result shows on the example of
\cref{thm:p15-ip-event-separation}.

\begin{hbcorollary}{thm:p15-w02b-p15-034-belief-not-prevision}{proved=annex:p15.w02b-p15-034.bel-prev}{Same belief function on events, different gamble prices}
On \(W=\{1,2,3\}\) let \(m(\{i\})=1/10\) for each \(i\),
\(m(\{i,j\})=1/5\) for each pair, and \(m(W)=1/10\). Then
\(\operatorname{Bel}_m\) is the common event lower probability of the two
sets \(C\) and \(D\) of \cref{thm:p15-ip-event-separation}, and for
\(f=(2,1,0)\),
\[
 \underline E_m(f)=1/2=\min_CE_pf<7/10=\min_DE_pf .
\]
So the lower envelope of \(D\) is a coherent lower prevision whose
restriction to events is the belief function \(\operatorname{Bel}_m\) but
which is not the gamble extension of that belief function.
\end{hbcorollary}

\begin{hbtheorem}{thm:p15-w02b-p15-034-strict-classes}{proved=annex:p15.w02b-p15-034.classes}{Necessity measures form a proper subclass of belief functions}
Every necessity measure \(\operatorname{Nec}_\pi\) on \(W\) is a belief
function \(\operatorname{Bel}_m\), for the nested mass assignment of
\cref{thm:p15-ip-possibility-bridge}. On \(W=\{a,b\}\) the belief function
with \(m(\{a\})=m(\{b\})=1/2\) equals no necessity measure.
Together with \cref{thm:p14-evidence-measure-separation} and
\cref{thm:p15-ip-event-separation} this gives three strict inclusions:
necessity measures among belief functions, belief functions among coherent
lower probabilities, and event restrictions that do not determine the
coherent lower prevision on gambles.
\end{hbtheorem}

\begin{hbdefinition}{def:p15-w02b-p15-034-dempster-conditioning}{Dempster conditioning}
Let \(m\) be a mass assignment on \(W\) and \(B\subseteq W\) with
\(\operatorname{Pl}_m(B)>0\). Let \(m^B\) place mass one on \(B\). The
\emph{Dempster conditional} of \(m\) on \(B\) is
\(m_B=m\otimes m^B\), the normalized combination of
\cref{def:p15-r47-combination}. Write
\(\operatorname{Bel}_m(\cdot\Vert B)=\operatorname{Bel}_{m_B}\) and
\(\operatorname{Pl}_m(\cdot\Vert B)=\operatorname{Pl}_{m_B}\).
\end{hbdefinition}

The double bar marks a different operation from the coherent conditional
\(\underline E(\cdot\mid B)\) of \cref{sec:p15-ip-conditioning}. Dempster's
combination treats the new information \(m^B\) as a second, independent
source and then discards the conflicting product mass; both are assumptions
of \cref{def:p15-r47-combination}, not consequences of coherence. The source
reports Walley's criticism that the conditioning rule derived from the
combination rule may fail to give conditional lower previsions coherent with
the unconditional ones, unless further conditional independence assumptions
are made \cite[§5.5, p.~35]{MirandaCoherentLowerPrevisionsCapture}. The
next result proves that failure on two explicit models, one with a
conditioning event of positive lower probability and one at the zero
boundary.

\begin{hbtheorem}{thm:p15-w02b-p15-034-dempster}{proved=annex:p15.w02b-p15-034.dempster}{Dempster conditioning is not the coherent update}
\begin{enumerate}
\item For \(\operatorname{Pl}_m(B)>0\) the combination in
\cref{def:p15-w02b-p15-034-dempster-conditioning} is defined, and for every
\(A\subseteq W\),
\[
 \operatorname{Pl}_m(A\Vert B)=\frac{\operatorname{Pl}_m(A\cap B)}{\operatorname{Pl}_m(B)},
 \qquad
 \operatorname{Bel}_m(A\Vert B)=1-\frac{\operatorname{Pl}_m(A^c\cap B)}{\operatorname{Pl}_m(B)}.
\]
If \(m\) is the nested mass assignment of a possibility distribution
\(\pi\), as in \cref{thm:p15-ip-possibility-bridge}, and
\(\operatorname{Poss}_\pi(B)>0\), then
\(\operatorname{Pl}_m(\cdot\Vert B)=\operatorname{Poss}_{\pi_B}\) and
\(\operatorname{Bel}_m(\cdot\Vert B)=\operatorname{Nec}_{\pi_B}\), where
\(\pi_B(w)=\pi(w)/\operatorname{Poss}_\pi(B)\) for \(w\in B\) and
\(\pi_B(w)=0\) otherwise.
\item On \(W=\{a,b,c\}\) let \(m(\{a,c\})=m(\{b\})=1/2\), \(B=\{a,b\}\)
and \(A=\{a\}\). Then
\(\operatorname{Bel}_m(A\Vert B)=\operatorname{Pl}_m(A\Vert B)=1/2\).
The natural extension \(\underline E_m\) has \(\underline E_m(B)=1/2>0\),
and its generalized Bayes update gives \(A\) conditional lower probability
\(0\) and conditional upper probability \(1/2\) on \(B\).
\item On \(W=\{a,b,c\}\) let \(\pi(a)=\pi(b)=1\), \(\pi(c)=1/2\), with
nested mass assignment \(m\), \(B=\{b,c\}\) and \(A=\{b\}\). Then
\(\operatorname{Bel}_m(A\Vert B)=1/2\) and
\(\operatorname{Pl}_m(A\Vert B)=1\). The natural extension of
\(\operatorname{Nec}_\pi\) gives \(B\) lower probability \(0\) and upper
probability \(1\), and both natural and regular conditioning give \(A\)
the conditional bounds \(0\) and \(1\) on \(B\).
\item In items~2 and~3, no separately coherent conditional lower
prevision on all gambles, jointly coherent with the natural extension in the
sense of \cref{def:p15-ip-conditional}, gives \(A\) the conditional lower
probability \(1/2\) on \(B\).
\end{enumerate}
\end{hbtheorem}

In item~2 the Dempster value is precise although the compatible models
disagree about the split of the mass on \(\{a,c\}\); it is the conditional
of the single compatible model that puts all of that mass on \(a\), the one
which maximizes the probability of \(B\). In item~3 the lower probability
of \(B\) is zero, the coherent conditional bounds are vacuous, and the
Dempster value \(1/2\) is more informative than any jointly coherent
single-event update allows, by \cref{thm:p15-ip-conditional-bounds}. The
possibility rule of item~1 is the ratio form
\(\operatorname{Poss}_\pi(A\cap B)/\operatorname{Poss}_\pi(B)\); a
possibility model updated by it is a different object from the coherent
update of the lower prevision that the same possibility distribution
generates. A user of either rule has to state which one is meant, because
the numbers in \([0,1]\) do not reveal it.

\section{A decision that depends on the independence judgement}
\label{sec:p15-w02b-p15-034-independence-decision}

Comparing two acts under an imprecise joint model requires the joint model,
and with imprecise marginals the joint model is not fixed by saying that two
variables are independent. The source distinguishes epistemic irrelevance,
epistemic independence, independence in the selection and strong
independence, and records that, when no conditioning event has probability
zero, each implies the one listed before it, all four coinciding for a
precise model \cite[§4, pp.~25--30]{MirandaCoherentLowerPrevisionsCapture}.
It also notes that sets of probabilities can leave two alternatives
unranked, and lists maximality and \(\Gamma\)-maximinity among the decision
criteria \cite[§6, p.~37]{MirandaCoherentLowerPrevisionsCapture}. The
example below uses the criteria of
\cref{def:p17-robust-dominance,thm:p17-robust-dominance}: the robust
comparison, under which \(f\) is strictly preferred to \(g\) when
\(\underline E(f-g)>0\), and the maxmin rule, which compares
\(\underline E(f)\) with \(\underline E(g)\).

\begin{hbdefinition}{def:p15-w02b-p15-034-two-binary}{Two imprecise binary marginals and three joint models}
Let \(X,Y\) take values in \(\{0,1\}\), and let them be logically
independent: \(\Omega=\{0,1\}^2\), every pair of values being possible.
Write a joint mass as \(p=(p_{00},p_{01},p_{10},p_{11})\), with
\(p_{xy}\) the probability of \(X=x,Y=y\), and write
\(p(X{=}x)\), \(p(Y{=}y)\) for its marginals. The marginal assessment for
each variable is the lower prevision
\[
 \underline P_1(h)=\min_{2/5\leq t\leq3/5}\bigl[t\,h(1)+(1-t)\,h(0)\bigr]
 \qquad(h\in\mathbb R^{\{0,1\}}),
\]
applied to gambles that depend on \(X\) alone, and separately to gambles
that depend on \(Y\) alone. The three joint models are the closed convex
sets
\[
\begin{aligned}
 \mathcal M_{\mathrm{mg}}&=\{p:\ 2/5\leq p(X{=}1)\leq3/5,\ 2/5\leq p(Y{=}1)\leq3/5\},\\
 \mathcal M_{\mathrm{ei}}&=\Bigl\{p:\ \tfrac25p(Y{=}y)\leq p_{1y}\leq\tfrac35p(Y{=}y)
   \text{ and}\\
 &\hphantom{{}=\Bigl\{p:\ }
   \tfrac25p(X{=}x)\leq p_{x1}\leq\tfrac35p(X{=}x)\ (x,y\in\{0,1\})\Bigr\},\\
 \mathcal M_{\mathrm{str}}&=\overline{\operatorname{conv}}
   \{p_X\otimes p_Y:\ p_X(1),p_Y(1)\in[2/5,3/5]\},
\end{aligned}
\]
the last being the product set of \cref{def:p17-mmeu-product} on this
finite space. Write \(\underline E_{\mathrm{mg}},\underline
E_{\mathrm{ei}},\underline E_{\mathrm{str}}\) for their lower envelopes
and \(\overline E_{\mathrm{mg}},\overline E_{\mathrm{ei}},\overline
E_{\mathrm{str}}\) for the conjugate upper envelopes. The acts are
\(f=I_{\{X=Y\}}\), a unit of utility if the two values agree, and the
constant \(g=9/20\).
\end{hbdefinition}

\begin{hbtheorem}{thm:p15-w02b-p15-034-independence-models}{proved=annex:p15.w02b-p15-034.ind-models}{What each independence judgement fixes}
In \cref{def:p15-w02b-p15-034-two-binary}:
\begin{enumerate}
\item \(\mathcal M_{\mathrm{str}}\subseteq\mathcal M_{\mathrm{ei}}
\subseteq\mathcal M_{\mathrm{mg}}\), and each of the three sets has the
marginal lower previsions \(\underline P_1\) for \(X\) and for \(Y\).
\item \(\mathcal M_{\mathrm{mg}}\) is the compatible set of the two marginal
assessments taken together, so \(\underline E_{\mathrm{mg}}\) is their
natural extension: the joint model with no judgement connecting \(X\) and
\(Y\).
\item Every closed convex \(\mathcal M\) with
\(\mathcal M_{\mathrm{str}}\subseteq\mathcal M\subseteq\mathcal M_{\mathrm{ei}}\)
gives each event \(\{X=x\}\) and \(\{Y=y\}\) lower probability \(2/5\);
its partition updates on \(\{\{Y=0\},\{Y=1\}\}\) and on
\(\{\{X=0\},\{X=1\}\}\) are unique and given by the generalized Bayes rule;
and under them \(X\) and \(Y\) are epistemically independent in the sense
of \cref{def:p15-ip-independence}.
\item Conversely, every closed convex set of joint masses with marginal lower
previsions \(\underline P_1\), under which \(X\) and \(Y\) are
epistemically independent with generalized-Bayes updates, is contained in
\(\mathcal M_{\mathrm{ei}}\). So \(\underline E_{\mathrm{ei}}\) is the
least informative epistemically independent model with these marginals.
\item The mass \(p^-=(1/5,3/10,3/10,1/5)\) is an extreme point of
\(\mathcal M_{\mathrm{ei}}\) and is not a product; hence
\(\mathcal M_{\mathrm{ei}}\ne\mathcal M_{\mathrm{str}}\).
\item The bounds on the agreement act are
\[
\begin{gathered}
 [\underline E_{\mathrm{str}}(f),\overline E_{\mathrm{str}}(f)]=[12/25,13/25],\qquad
 [\underline E_{\mathrm{ei}}(f),\overline E_{\mathrm{ei}}(f)]=[2/5,3/5],\\
 [\underline E_{\mathrm{mg}}(f),\overline E_{\mathrm{mg}}(f)]=[0,1].
\end{gathered}
\]
\end{enumerate}
\end{hbtheorem}

\begin{hbtheorem}{thm:p15-w02b-p15-034-independence-decision}{proved=annex:p15.w02b-p15-034.ind-decision}{The comparison of the agreement bet with a sure payment}
In \cref{def:p15-w02b-p15-034-two-binary}:
\begin{enumerate}
\item Under \(\mathcal M_{\mathrm{str}}\),
\(\underline E_{\mathrm{str}}(f-g)=3/100>0\): every retained model gives
\(f\) the higher expectation, and the maxmin rule also chooses \(f\).
\item Under \(\mathcal M_{\mathrm{ei}}\),
\(\underline E_{\mathrm{ei}}(f-g)=-1/20<0<3/20=\overline E_{\mathrm{ei}}(f-g)\):
some retained model ranks \(f\) strictly above \(g\) and another strictly
below, so neither act robustly dominates the other. The maxmin rule chooses
\(g\).
\item Under \(\mathcal M_{\mathrm{mg}}\),
\(\underline E_{\mathrm{mg}}(f-g)=-9/20\) and
\(\overline E_{\mathrm{mg}}(f-g)=11/20\), with the same two conclusions as
in item~2.
\item Every precise product model in \(\mathcal M_{\mathrm{str}}\) chooses
\(f\); the precise model \(p^-\) chooses \(g\).
\end{enumerate}
\end{hbtheorem}

The robust preference for the agreement bet therefore needs the strong
product judgement: that every combination of compatible marginal models is
a compatible joint model, and nothing else is. Epistemic independence does
not supply it. \(\mathcal M_{\mathrm{ei}}\) satisfies epistemic
independence and leaves the two acts unranked, while
\(\mathcal M_{\mathrm{str}}\) also satisfies it and ranks them. The
unranked pair under \(\mathcal M_{\mathrm{ei}}\) and
\(\mathcal M_{\mathrm{mg}}\) is retained imprecision: the expected advantage
of \(f\) over \(g\) has lower bound below zero and upper bound above it,
and the model does not say that the acts are equally good. A choice still
has to be made, and the maxmin rule makes one, but that choice is a
property of the added rule, not of the model. Indifference would be a
different judgement, a precise expected advantage of zero, which the
interval does not assert.

The comparison with precise probability is item~4. A precise model under
which \(X\) and \(Y\) are judged independent is a product, and every
product with these marginals chooses \(f\). With imprecise marginals the
notions separate: \(p^-\) has conditional probabilities of \(X=1\) given
\(Y=y\) equal to \(2/5\) and \(3/5\), each inside the marginal interval, so
it belongs to an epistemically independent set, although as a single
precise model it makes \(X\) and \(Y\) dependent. Logical independence is
used throughout; it is a separate assumption from any of these, and it is
what leaves all four joint values free to vary. The example is constructed
for this chapter; it is not taken from the source.

\paragraph{Relationships.}
\Cref{thm:p15-w02b-p15-034-necessity} is the partition form of the
single-event identity in \cref{thm:p15-ip-conditional-bounds}: there the
sold and bought transactions involve one event, here the conditional
transaction may be active on every cell. The generalized-Bayes root used in
\(\underline P_{\mathcal B}\) is the root analysed in
\cref{thm:p15-ip-generalized-bayes} for one event. The finitely additive
component \(\nu\) in \cref{thm:p15-w02b-p15-034-counterexample} is of the
kind allowed by \cref{def:p15-ip-compatible}, where no countable additivity
is imposed.
\Cref{thm:p15-w02b-p15-034-belief-extension} gives the gamble form of the
allocated envelopes of \cref{thm:p14-evidence-finite-measures}, and
agrees with the convex-core identity of \cref{thm:p17-rat84-choquet} for the
convex capacity \(\operatorname{Bel}_m\). The second inclusion of
\cref{thm:p15-w02b-p15-034-strict-classes} has an inequality form: every
lower envelope is superadditive on pairs of disjoint events
(\cref{thm:p15-w02b-p15-032-disjoint}), while belief functions satisfy the
inclusion--exclusion inequality at every order, and the order-two
inequality already fails for some lower envelopes.
\Cref{thm:p15-w02b-p15-034-dempster} uses the combination rule of
\cref{def:p15-r47-combination} for conditioning only; whether two sources
justify that rule's independence premise is the separate question treated
with it. The independence notions of \cref{def:p15-ip-independence} are
compared at a null event in \cref{thm:p15-ip-independence-boundary}; the
models of \cref{def:p15-w02b-p15-034-two-binary} have no null conditioning
event, so the separation found there between epistemic and strong
independence is not a zero-probability effect. The product set
\(\mathcal M_{\mathrm{str}}\) is the finite case of
\cref{thm:p17-mmeu-product-priors}. The robust comparison and the maxmin
rule used in \cref{thm:p15-w02b-p15-034-independence-decision} are the
difference and maxmin comparisons of \cref{def:p15-w02b-p15-032-rules};
\cref{thm:p15-w02b-p15-032-ellsberg} shows that the first satisfies the
sure-thing pasting condition and the second reproduces Ellsberg's
three-colour pattern. The example here adds a different dependence: with
the marginal assessments held fixed, whether the robust comparison ranks
the pair at all turns on the independence judgement used to build the
joint model.
